\subsection{All New Requirements of the Saved PC2 Policy}\label{ta:pc2_all_act}
The previous history analysis fixes one requirement to distinguish A from H. A separate inspection of the same saved policy establishes the active-lifetime obligation for all thirteen new requirements. It introduces no solver run, input change or timing trial.

The only new-start edges are the chain $q_{20}\to q_{19}\to\cdots\to q_7$, in this order: stamping avoidance for arms 1 and 2; calibration availability; clean-once for arms 1 and 2; drill-once for arms 1 and 2; output-if-finished for arms 1 and 2; new tool order for arms 1 and 2; and paint-once for arms 1 and 2. The selected start above is the eighth edge. Every state in this chain has two NEW/ARM components, all 24 ordinary observers zero and rank $q-7$ at $q$. Each newly installed monitor is 0 and stays 0. There is no ordinary event, transfer or UC edge after the first new-start.

\begin{proposition}[All PC2 active obligations]\label{prop:pc2-all-act}
Every maximal run of the saved PC2-Rolling policy from any of its 126 entries activates all thirteen new requirements and satisfies each source invariant from its own activation through handover. The matching verified new endpoint preserves these obligations afterward.
\end{proposition}
\begin{proof}
The saved domain is entry-reachable and closed, every edge strictly decreases a natural rank, and no non-goal deadlocks. Hence every maximal run reaches the unique goal $q_7$. Every $q_j$, $7\le j\le19$, has only its preceding chain edge as an incoming policy edge, and no entry lies on the chain. Thus every run traverses all thirteen starts.

For each arm, the source invariants forbid stamping; forbid cleaning, drilling or painting an already completed workpiece; require all three completed-tool fluents at output; and constrain tool-pending fluents by the new tool order. At activation all completed-tool and pending fluents are false, and no ordinary action occurs on a start label. Each invariant therefore holds. The calibration invariant is $\Box(\neg Calibrating_1\lor\neg Calibrating_2)$. Its source fluents start false and are set only by ordinary calibration requests. The old plant has no such request, so every old entry retains false values; none occurs in this policy either, while completion clears them. They too are false on the start chain. Start commands preserve these fluents and produce no ordinary action, so all active source invariants continue to hold until $q_7$. This goal matches \texttt{new-00000009}, with both components, all ordinary observers and all thirteen new monitors aligned. The supplied verified endpoint and continuation preservation give the post-handover part.
\end{proof}

This covers all new requirements of this one saved policy at their reached activation tuples. It does not prove initializer exactness over every permitted tuple, safe pre-activation histories, all 138 NEW occurrences, compiler correctness or runtime conformance. The 39/87 history partition above is unchanged. PC2 completes calibration before activating any new requirement; its active update suffix contains only starts. The GSM construction below complements it with ordinary behavior under an active new requirement.

The earlier violations have a concrete source distinction: old PC2 requires drilling, polishing and cleaning, whereas the new output rule requires drilling, cleaning and painting. The update contract preserves its interval calibration obligation and whichever old requirements remain active; it begins each new duty at its start. This defines the specified obligation, without establishing that an operator would choose that scope.
